\section{System Design and Methodology}
\label{sec:system-design}

PyTrade separates two problems that are usually learned jointly: extracting a temporal trading signal from a single instrument, and allocating capital across many instruments. Section~\ref{sec:overview} introduces the system's design philosophy and the formal decision-process framework shared by both layers. Sections~\ref{sec:state} and~\ref{sec:reward} describe the state-construction and reward machinery both layers draw on before either layer is described in detail. Section~\ref{sec:saa} then gives a complete, self-contained account of the temporal layer, the Single-Asset Agent (SAA); Section~\ref{sec:paa} gives the same account for the cross-sectional layer, the Portfolio Allocator Agent (PAA), including the exact mechanism by which the SAA's signal is produced and injected into it. Sections~\ref{sec:simulator}--\ref{sec:universe} close the chapter with the simulator, cost model, and asset universe that make the whole system, both layers together, trainable and testable end to end.

\subsection{System Overview}
\label{sec:overview}

PyTrade decomposes portfolio allocation into two learned layers that are trained, and are best understood, separately. Layer 1, the Single-Asset Agent (SAA), is a recurrent policy that learns to trade one instrument in isolation: given only that instrument's own price history and its own current position, it decides whether to buy, hold, or sell. Layer 2, the Portfolio Allocator Agent (PAA), is a self-attention policy that observes the entire eleven-asset universe at once and decides how to split capital across all of it. Crucially, the PAA does not learn its own opinion about each asset's timing from scratch: it consumes the SAA's per-asset trading signal as one of its inputs, alongside raw market data, and learns how much to trust, discount, or override that signal for each asset given the current state of the whole universe. This chapter follows that hierarchy in the order it is actually built: this section finishes with the mathematical framework shared by both layers, Sections~\ref{sec:state} and~\ref{sec:reward} describe the state-construction and reward machinery both layers draw on, Section~\ref{sec:saa} gives a complete account of the SAA on its own, Section~\ref{sec:paa} gives the same account for the PAA, and Sections~\ref{sec:simulator}--\ref{sec:universe} describe the infrastructure that trains and evaluates both together.

Both layers are trained inside the same kind of environment, so it is worth fixing the vocabulary once before any formalism. The environment advances in discrete steps, each corresponding to one trading day. At the start of a step the environment hands the agent an \emph{observation}: a fixed-length numeric vector describing recent market behavior and the agent's own current holdings. The agent's network maps this observation to an \emph{action}: for the SAA a single number describing how much to buy or sell of its one asset, for the PAA a vector describing how to reweight the entire portfolio. The environment then executes the implied trade against that day's price, subject to the transaction-cost model of Section~\ref{sec:friction}, advances the calendar by one day, revalues every position at the new closing price, and computes a scalar \emph{reward} that scores how well that single step went (Section~\ref{sec:reward}). This observation-action-reward cycle repeats for a fixed number of steps, an \emph{episode} (252 trading days, roughly one calendar year, Section~\ref{sec:simulator}), sampled from a contiguous block of historical dates; many episodes, sampled from different periods of the 2000--2025 history, are played out in parallel, and Proximal Policy Optimization (PPO), the policy-gradient algorithm introduced in the literature review, adjusts the network's parameters after each batch of collected experience to increase the expected sum of rewards it receives. One detail is easy to miss and important to state explicitly: the historical price series itself is exogenous data, replayed unchanged episode after episode, so nothing the agent does changes tomorrow's closing price. The environment's transition dynamics are therefore deterministic given the historical data and the agent's own trade, not a market-impact simulation of the kind that would let the agent's own order flow move the market; only the agent's own cash, positions, and (for the SAA) internal recurrent state evolve as a genuine consequence of its actions.

Splitting the temporal problem from the cross-sectional problem is a direct architectural response to two very different structures in the underlying decision problem, not merely a modeling convenience. Learning a good trading rule for one asset from its own price history has a strong sequential-memory requirement: an unusual move today only makes sense read against the last weeks of context, which recurrent networks such as the Long Short-Term Memory (LSTM) cell used here are specifically built to carry forward. Learning how eleven assets should be weighted \emph{against each other}, by contrast, is at every instant a set problem rather than a sequence problem: given today's picture of all eleven assets, the right relative allocation does not fundamentally depend on the order in which the assets happen to be listed, which is exactly the permutation-respecting structure a self-attention mechanism (Section~\ref{sec:paa}) provides and a recurrent network does not. Asking one, monolithic network to learn both structures at once would force it to discover this decomposition on its own, from a reward signal alone; PyTrade instead builds the decomposition directly into the architecture, training a temporal specialist first and a cross-sectional specialist second, and only asking the second network to learn the comparatively narrower problem of how much to trust the first.

We formalize both layers' decision problems with the same general tuple and instantiate it twice, with agent-specific state, action, and observation spaces, in Sections~\ref{sec:saa} and~\ref{sec:paa}. A Partially Observable Markov Decision Process (POMDP) is a tuple

\begin{equation}
\mathcal{M} = (\mathcal{S}, \mathcal{A}, \mathcal{O}, \mathcal{P}, \mathcal{R}, \mathcal{Z}),
\end{equation}

where $\mathcal{S}$ is the true simulator state, every quantity the environment itself tracks at a given step whether or not the agent ever sees it; $\mathcal{A}$ is the action space, the set of decisions the policy is allowed to make; $\mathcal{O}$ is the observation space, the (generally smaller) space of vectors actually handed to the network; $\mathcal{P}$ is the transition function, deterministic throughout this system given the historical price data and the agent's own action, as just discussed; $\mathcal{R}$ is the reward function (Section~\ref{sec:reward}); and $\mathcal{Z}: \mathcal{S} \to \mathcal{O}$ is the deterministic observation map that decides which parts of the true state the agent is allowed to see. Both agents retain PPO's standard discount factor $\gamma$ and Generalized Advantage Estimation's (GAE) decay parameter $\lambda$ (Section~\ref{sec:saa}) rather than setting $\gamma=1$: each experiment configures its own $\gamma \in [0.98, 0.998]$ and $\lambda \in [0.95, 0.98]$, so later steps within an episode still contribute to the advantage estimate with exponentially decaying weight. This is a deliberate variance-reduction choice, not a consequence of the finite, 252-step episode length: a finite horizon would in principle permit an undiscounted ($\gamma=1$) objective, but this thesis does not take advantage of that possibility.


Partial observability in this system has two distinct sources, and it is worth being precise about which one the POMDP formalism above actually captures. The dominant source is not a modeling choice at all: daily-frequency financial markets are influenced by an effectively unbounded set of macroeconomic, political, and microstructural factors that are never fully measurable, and even a perfectly informed observer could not forecast their future evolution from present information. No feature engineering closes that gap. A second, compounding source is a deliberate engineering choice: this system reduces every asset to a small, hand-selected panel of stationary technical features (Section~\ref{sec:state}), which discards information a full order book, news feed, or fundamentals dataset would carry. Both of these sources mean $\mathcal{O}$ is, by construction, a lossy view of whatever true generative process drives asset prices, independent of any network architecture. What the POMDP tuple above formalizes on top of that is a third, narrower, and purely technical source: the SAA's LSTM cell state is an internal summary of past observations that is never re-exposed to the environment or the reward function, so from the environment's point of view the SAA's policy is not a memoryless function of the current observation alone. This third source is real and matters for the transfer argument in Section~\ref{sec:saa-stability}, but it should not be read as a claim that the agents come close to seeing the true state of the market; they do not, and no realistic system trained on daily OHLCV-derived features would. The PAA's own policy, by contrast, \emph{is} a memoryless function of its current observation (Section~\ref{sec:paa}); any memory that reaches it about the past arrives indirectly, already compressed into the scalar signal produced by a frozen, recurrent SAA copy (Section~\ref{sec:info-flow}).

\subsection{State Representation: A Shared Feature Engineering Pipeline}
\label{sec:state}

Raw OHLCV (open, high, low, close, volume) data is not stationary, and stacking more raw history does not fix that; it just gives the network more non-stationary numbers to memorize. A stochastic process is (weakly) \emph{stationary} if its mean and variance do not drift over time and its autocovariance structure depends only on the lag between two observations, never on calendar time itself. A raw price series violates this, since its mean is, roughly speaking, wherever the trend has taken it throughout the decade. This matters for a neural network for a very concrete reason: a feature whose typical scale keeps shifting between the period it was trained on and the period it is evaluated on behaves, from the network's perspective, like a moving target, and the resulting distribution shift is exactly what the purged, block-level evaluation of Section~\ref{sec:simulator} is designed to expose rather than paper over. Every feature described in this section is produced once, offline, by a dedicated enrichment pipeline (\texttt{DataEnricher} in \texttt{data\_processor\_2.py}) and cached by the \texttt{MarketDataCache} of Section~\ref{sec:simulator}; the environment itself never computes a market feature at runtime. Portfolio features such as cash levels, asset weights or transaction costs are calculated per step at runtime.  

Two generic transformations recur throughout the pipeline and are worth fixing once, with notation used consistently for the rest of this chapter. The first is the log return over a horizon of $N$ days, $r_t^{(N)} = \log(P_t/P_{t-N})$, used instead of a simple percentage return because log returns compound additively across time (a $+10\%$ day followed by a $-10\%$ day very nearly cancels out under log returns, but leaves a small net loss under simple returns) and treat gains and losses of equal relative magnitude symmetrically. The second is rolling $z$-score standardization,
\begin{equation}
\hat x_t = \frac{x_t - \mu_w(x)_t}{\sigma_w(x)_t + \epsilon},
\end{equation}
where $\mu_w(\cdot)_t$ and $\sigma_w(\cdot)_t$ denote the rolling mean and standard deviation of a series over the trailing $w$ days and $\epsilon$ guards against division by zero; nearly every feature below is, in its final form, some raw quantity passed through this standardization, and several are further squashed with $\tanh$ into a bounded $(-1,1)$ range. Nearly all rolling statistics in the pipeline share a single window length, $w_{\text{norm}} = \lfloor 252/7 \rfloor = 36$ trading days (roughly one and a half months), which keeps the question ``after how much history is this $z$-score averaging over?'' answerable with one number throughout this section; the exact train, validation, and test partitioning this window length has to respect is elaborated in Section~\ref{sec:simulator}. Every feature is computed independently per asset, so the same feature vector is meaningful across the eleven-asset universe (Section~\ref{sec:universe}); which features actually reach the SAA and which reach each of the PAA's two token types is decided entirely by a boolean map in that experiment's JSON configuration, as detailed for each agent in Sections~\ref{sec:saa} and~\ref{sec:paa}.

\subsubsection{Multi-Horizon Return, Volume, and Volatility Features}
\label{sec:state-returns}

A common feature family is computed at three horizons $N \in \{1, 5, 21\}$ trading days (the enrichment pipeline itself also computes $N \in \{3, 10, 20, 40\}$, but only the 1/5/21-day variants are enabled in the active feature configuration, matching daily, weekly, and monthly trading rhythms). For each horizon:

\begin{equation}
r_t^{(N)} = \log\frac{P_t}{P_{t-N}}, \qquad r_t^{(N),\text{SPY}} = \log\frac{P_t^{\text{SPY}}}{P_{t-N}^{\text{SPY}}}, \qquad \tilde{r}_t^{(N)} = r_t^{(N)} - r_t^{(N),\text{SPY}},
\end{equation}

the asset's own log return, SPY's log return over the same horizon, and the asset's log return relative to SPY. The market-relative return $\tilde{r}_t^{(N)}$, not the raw return $r_t^{(N)}$, is what actually enters the feature set (as \texttt{log\_close\_return\_Nd} is joined by its z-scored relative counterpart):

\begin{equation}
z\text{\_rel\_log\_close}_t^{(N)} = \text{clip}\!\left(\frac{\tilde{r}_t^{(N)} - \mu_{36}(\tilde{r}^{(N)})_t}{\sigma_{36}(\tilde{r}^{(N)})_t + \epsilon}, \, -4, 4\right),
\end{equation}

where $\mu_{36}(\cdot)_t$ and $\sigma_{36}(\cdot)_t$ denote the rolling mean and standard deviation over the trailing $w_{\text{norm}}=36$ days. Removing SPY's return before $z$-scoring is deliberate: it prevents a broad market rally or sell-off from being read by the network as an asset-specific signal, which matters directly for this thesis's modular design, since genuine cross-asset structure should come from the PAA's attention mechanism (Section~\ref{sec:paa}), not leak in through a feature that is secretly just ``the market went up today.''

A volatility-breakout companion feature, \texttt{z\_vol\_breakout\_Nd}, flags when short-term realized volatility of the relative return has broken out of its recent regime:

\begin{equation}
z\text{\_vol\_breakout}_t^{(N)} = \text{clip}\!\left(\frac{\log\big(\sigma_5(\tilde{r}^{(N)})_t + \epsilon\big) - \mu_{36}(\tilde{r}^{(N)})_t}{\sigma_{36}(\tilde{r}^{(N)})_t + \epsilon}, \, -4, 4\right),
\end{equation}

comparing the log of a 5-day realized-volatility estimate against the 36-day rolling mean and standard deviation of the relative return series itself. An analogous pair of features is computed for traded volume: $\text{log\_volume\_intensity}_t^{(N)} = \log\!\big(V_t / (\mu_{36}(V)_t + \epsilon)\big)$ measures today's volume against its own trailing average, capturing the intuition that a genuine regime change is more often confirmed by unusual volume than by price movement alone; it is again made market-relative by subtracting SPY's volume intensity, $z$-scored over 36 days, and squashed with $\tanh$ into \texttt{norm\_volume\_feature\_Nd} $\in (-1,1)$.

\subsubsection{Intraday Shape and Classical Technical Indicators}

A second group of features is computed once per day (no horizon parameter) from that day's own OHLC bar, capturing information a close-only feature discards entirely: where within the day's range the close occurred, and whether the day opened with a gap.

\begin{equation}
\text{intra\_vol}_t = \log\frac{H_t}{L_t}, \qquad
\text{cl\_strength}_t = \frac{C_t - L_t}{H_t - L_t}, \qquad
\text{gap}_t = \log\frac{O_t}{C_{t-1}}, \qquad
\text{shadow}_t = \frac{H_t - \max(O_t, C_t)}{H_t - L_t},
\end{equation}

the log high-low range, the fraction of the day's range the close sits above the low (bounded in $[0,1]$), the overnight log return from the previous close to today's open, and the upper-wick fraction of the day's range. $\text{intra\_vol}_t$ and $\text{gap}_t$ are $z$-scored over the 36-day window and passed through $\tanh$ (\texttt{z\_log\_intra\_vol}, \texttt{z\_log\_overnight\_gap}); $\text{cl\_strength}_t$ and $\text{shadow}_t$ are already bounded in $[0,1]$ and are simply rescaled to $[-1,1]$ via $2x-1$ (\texttt{norm\_closing\_strength}, \texttt{norm\_shadow\_ratio}).

Three classical technical indicators are computed alongside these bar-shape features, standard tools in discretionary trading practice that are worth defining explicitly since they are rarely discussed in the reinforcement-learning literature, and each rescaled here to be dimensionless and scale-free across the eleven-asset universe.

The \emph{Relative Strength Index} (RSI) is a momentum oscillator that compares the average size of recent up days to the average size of recent down days over a lookback window, here 14 trading days: writing $\overline{\text{gain}}_{14}$ and $\overline{\text{loss}}_{14}$ for those two rolling averages of positive and negative daily price changes respectively, $\text{RS} = \overline{\text{gain}}_{14}/\overline{\text{loss}}_{14}$ and $\text{RSI} = 100 - 100/(1+\text{RS})$, which is bounded in $[0,100]$ by construction; values above 70 are conventionally read as an asset being ``overbought,'' below 30 as ``oversold.'' The pipeline rescales and smooths this into $\texttt{rsi\_14} = \tfrac{1}{3}\sum\big((\text{RSI}_{14}-50)/50\big)$, a 3-day moving average of the recentered, $[-1,1]$-bounded oscillator, matching the scale most other features in this pipeline share.

A \emph{Bollinger Band} places an envelope around a rolling mean price at a fixed multiple $k$ of the rolling standard deviation, $\mu_w(P)_t \pm k\,\sigma_w(P)_t$; discretionary trading practice typically watches where the current price sits inside this envelope. This pipeline does not feed the band boundaries or the price's position within them; instead it feeds the band's own \emph{width} relative to the price level, $4\sigma_{36}(P)_t / \mu_{36}(P)_t$ (with $k=2$, so the full envelope spans $4\sigma$), itself $z$-scored over the same 36-day window and clipped to $[-3,3]$ as \texttt{bb\_width\_reworked}. A widening band signals an expanding volatility regime; a narrowing band signals compression, which discretionary practice associates with a regime change about to break in either direction. Feeding the normalized width, rather than the raw envelope or the price's position in it, keeps the feature meaningful whether an instrument trades at \$20 or \$2{,}000 a share, and keeps it symmetric with respect to whichever direction any eventual breakout takes.

The \emph{Average True Range} (ATR) is a volatility-magnitude indicator that averages, over a rolling window, the \emph{true range} $\max\big(H_t - L_t,\ |H_t - C_{t-1}|,\ |L_t - C_{t-1}|\big)$ rather than the plain high-low range $H_t - L_t$; the extra two terms ensure that a session which gaps sharply away from the previous close is still counted as volatile even when that day's own high-low range happens to be narrow. The resulting series is $z$-scored over 36 days and clipped to $[-3,3]$ into \texttt{atr\_z\_norm}. All three indicators are included for the same underlying reason: each summarizes, in one scalar, a regime property (momentum direction and strength for RSI, price-relative volatility for the Bollinger width, gap-aware volatility magnitude for the ATR) that plain multi-horizon returns and volume do not directly expose to the network, and that decades of discretionary practice suggest is informative on its own.

\subsubsection{Cross-Asset Features}

Three features make each asset's representation depend explicitly on the rest of the universe, giving the PAA's cross-sectional attention (Section~\ref{sec:paa}) genuine correlation structure to work with even before any learned attention weights are applied.

The first is a rolling 21-day Pearson correlation, the standard linear co-movement statistic 
\begin{equation}
    \rho = \mathrm{Cov}(r^{(1)},r^{(1),\text{SPY}})/\big(\sigma(r^{(1)})\,\sigma(r^{(1),\text{SPY}})\big),
\end{equation}
between the asset's daily log return and SPY's. Correlation coefficients are bounded in $(-1,1)$ and their sampling distribution is skewed near those bounds, which makes rolling $z$-scoring of the raw coefficient statistically awkward; the \emph{Fisher $z$-transform} addresses this directly,

\begin{equation}
\rho_t = \text{clip}\big(\text{Corr}_{21}(r^{(1)}, r^{(1),\text{SPY}})_t,\, -0.99, 0.99\big), \qquad
z\text{\_fisher\_corr\_spy}_t = \tfrac{1}{2}\log\frac{1+\rho_t}{1-\rho_t},
\end{equation}

which maps the bounded correlation onto an unbounded, approximately variance-stabilized scale; this is what makes it safe to treat like any other roughly Gaussian feature within this pipeline's shared $z$-score/$\tanh$ normalization scheme, rather than a bounded ratio whose variance changes depending on how close it already sits to $\pm 1$.

The second is a rolling 36-day \emph{beta} against SPY, computed from the same covariance/variance ratio a Capital Asset Pricing Model (CAPM) regression would use, $\beta = \mathrm{Cov}(r^{(1)}, r^{(1),\text{SPY}}) / \mathrm{Var}(r^{(1),\text{SPY}})$, then $z$-scored over 36 days into \texttt{z\_beta\_spy}; a beta above one indicates the asset has historically amplified SPY's moves, below one that it has dampened them.

The third is a cross-sectional \emph{dispersion} feature, $\log\big(\text{std}_i(r_{i,t}^{(1)}) + \epsilon\big)$, the standard deviation, taken across all eleven assets on date $t$, of that day's log return, itself $z$-scored over 36 days into \texttt{z\_market\_dispersion}. This gives every asset a shared, date-aligned read of ``how dispersed is the whole universe's behavior today'': a low value means most assets are moving together (a broad macro regime dominating), a high value means today is a stock-picker's day where individual asset behavior diverges, context the PAA's attention weights can use directly rather than having to re-derive it from eleven separate return series on their own.

\subsubsection{Memory-Preserving Stationary Features via Fractional Differentiation}
\label{sec:state-ffd}

Standard integer differencing, i.e., using $\log P_t - \log P_{t-1}$ as the only price-derived feature, achieves stationarity but destroys long-range memory: a return series has, by construction, no information left about the price level or slow trend that produced it. López de Prado \cite{LopezdePrado.2018} shows this is not a necessary trade-off. Treating the difference order as a continuous parameter $d \in \mathbb{R}$ rather than the integer $d=1$ of ordinary returns, the fractional difference operator is defined through the binomial series expansion of $(1-L)^d$ in the backshift operator $L$ (where $L X_t = X_{t-1}$):

\begin{equation}
(1-L)^d = \sum_{k=0}^{\infty} \binom{d}{k}(-L)^k = \sum_{k=0}^{\infty} w_k L^k, \qquad
w_0 = 1, \quad w_k = -w_{k-1}\,\frac{d - k + 1}{k}.
\end{equation}

For $d=1$ this collapses to $w_0=1, w_1=-1$, i.e., ordinary differencing; for non-integer $d \in (0,1)$, the weights $w_k$ decay slowly toward zero rather than vanishing after one lag, so the transformed series retains a fading memory of every past price rather than none of it.

Whether a given choice of $d$ has actually achieved stationarity is an empirical question, and this pipeline answers it with the \emph{Augmented Dickey-Fuller} (ADF) test, never by visual inspection of a chart, every time this subsection states a $p$-value. The ADF test asks a narrow, well-posed question: does a series behave like a \emph{unit-root} process, $y_t = y_{t-1} + \varepsilon_t$, a random walk that never reverts to a fixed level and whose variance grows without bound over time, or does it revert toward its own mean? Concretely, the test regresses the series' first difference on its own lagged level plus a number of further lagged differences,
\begin{equation}
\Delta y_t = \phi\, y_{t-1} + \sum_{i=1}^{p} \psi_i\, \Delta y_{t-i} + \varepsilon_t,
\end{equation}
where the number of lag terms $p$ is chosen automatically by minimizing the Akaike Information Criterion (\texttt{statsmodels.adfuller(..., autolag='AIC')}), which trades off how well the regression fits against how many nuisance lag terms are added purely to remove residual autocorrelation from $\varepsilon_t$. The null hypothesis is $\phi = 0$, a unit root is present and the series is non-stationary; the alternative is $\phi < 0$, the series pulls back toward its own mean and is stationary. Because $\phi$'s test statistic does not follow a standard $t$-distribution under this null, the test compares it against its own tabulated asymptotic distribution and reports a $p$-value; this pipeline rejects the unit-root null, and accepts a series as stationary, whenever $p < 0.05$.

Because the exact weight series in the operator above is infinite, \texttt{DataEnricher.\_get\_global\_ffd\_weights} uses López de Prado's fixed-width window approximation: weights are generated by the recursion above until either they fall below a threshold or a hard cap is reached,

\begin{equation}
\ell^{\ast} = \min\Big(\big\{k : |w_k| < \tau\big\} \cup \{k_{\max}\}\Big), \qquad \tau = 10^{-3},\ \ k_{\max}=120,
\end{equation}

and the fractionally differentiated series is then the finite, fixed-width convolution

\begin{equation}
\widetilde{X}_t = \sum_{k=0}^{\ell^{\ast}} w_k \,\log P_{t-k},
\end{equation}

applied to the log price rather than the raw price (log-prices keep the weights scale-free across assets with very different price levels). This is a genuine memory/stationarity trade-off, not a free lunch: smaller $d$ preserves more memory but risks failing the ADF test; larger $d$ is safely stationary but approaches ordinary returns and their short memory. The pipeline therefore searches for the smallest $d$ that empirically passes the ADF test, the ``golden $d$,'' rather than fixing $d$ by convention. Concretely, \texttt{find\_golden\_d} and \texttt{find\_multi\_asset\_golden\_d} execute the following procedure, shown end to end in Figure~\ref{fig:ffd-pipeline}:

\begin{enumerate}
\item Fix a candidate order $d$, starting at $0.45$ and stepping by $0.05$ up to $0.8$, and generate its weight vector $w_k$ via the recursion above, truncated at $\ell^\ast$.
\item Convolve $w_k$ with the asset's log-price series to obtain the candidate series $\widetilde X_t$.
\item Run the ADF test on $\widetilde X_t$: if $p<0.05$, accept $d$ as this asset's golden $d$ and stop; otherwise increase $d$ by the step size and repeat.
\item Repeat steps 1--3 independently for three structurally different anchor assets, a broad equity index (SPY), a precious metal (Gold), and a commodity (Crude), so the search is not biased toward any one asset's own idiosyncratic memory.
\item Adopt $d^{\ast} = \max$ of the three anchors' golden $d$'s as a single, shared differencing order for the whole universe.
\item Generate the weight vector once, at $d^{\ast}$, and apply it uniformly to every asset in the eleven-asset universe, including the eight not used in the search.
\item Normalize the resulting \texttt{ffd\_close} series exactly like the other rolling features, a 36-day $z$-score followed by $\tanh$, into \texttt{norm\_ffd\_feature}.
\item Re-run the ADF test on \texttt{norm\_ffd\_feature} independently for all eleven assets as a closing correctness gate, requiring a $100\%$ pass rate at $p<0.05$; if even one asset fails, the pipeline raises rather than silently shipping a non-stationary feature into training.
\end{enumerate}

Step 5's choice of the maximum, rather than tuning $d$ per asset, is deliberate on two counts: tuning per asset would make the resulting feature's memory horizon inconsistent across the universe, undermining any claim that the PAA's cross-sectional attention (Section~\ref{sec:paa}) is comparing like with like across assets; and since a larger $d$ is always at least as safely stationary as a smaller one, taking the maximum guarantees the shared weight vector already passes for the three anchors, while step 8 confirms empirically, rather than merely assumes, that it also transfers to the eight assets never used in the search.

\begin{figure}[htbp]
\centering
\begin{tikzpicture}[
    node distance=6mm and 8mm,
    box/.style={draw, rounded corners, fill=black!5, minimum height=9mm, minimum width=64mm, align=center, font=\small},
    small/.style={draw, rounded corners, fill=black!5, minimum height=8mm, minimum width=28mm, align=center, font=\footnotesize},
    dec/.style={draw, rounded corners, fill=black!10, minimum height=9mm, minimum width=64mm, align=center, font=\small},
    arr/.style={-{Latex[length=2mm]}}
]
\node[box] (logp) {Log price $\log P_t$, all 11 assets};
\node[box, below=of logp] (weights) {Binomial weights $w_k$ via recursion\\ $w_0{=}1,\ w_k=-w_{k-1}\frac{d-k+1}{k}$, cut at $\lvert w_k\rvert{<}\tau$};
\node[small, below=10mm of weights] (gold) {Gold};
\node[small, left=6mm of gold] (spy) {SPY};
\node[small, right=6mm of gold] (crude) {Crude};
\node[dec, below=10mm of gold] (search) {Golden-$d$ search: smallest $d\in[0.45,0.8]$ (step 0.05)\\ passing an ADF stationarity test ($p{<}0.05$), per anchor asset};
\node[box, below=of search] (maxd) {$d^{\ast} = \max(d_{\text{SPY}}, d_{\text{Gold}}, d_{\text{Crude}})$\\ one shared weight vector for the whole universe};
\node[box, below=of maxd] (conv) {Fixed-width convolution $\widetilde{X}_t = \sum_{k=0}^{\ell^\ast} w_k \log P_{t-k}$\\ applied per asset with $d^\ast$};
\node[box, below=of conv] (norm) {Rolling 36-day $z$-score, then $\tanh$\\ $\to$ \texttt{norm\_ffd\_feature}};
\node[dec, below=of norm] (adf) {Per-asset ADF verification on \texttt{norm\_ffd\_feature}\\ requires 100\% pass rate ($p{<}0.05$), else pipeline aborts};
\draw[arr] (logp) -- (weights);
\draw[arr] (weights) -- (spy);
\draw[arr] (weights) -- (gold);
\draw[arr] (weights) -- (crude);
\draw[arr] (spy) -- (search);
\draw[arr] (gold) -- (search);
\draw[arr] (crude) -- (search);
\draw[arr] (search) -- (maxd);
\draw[arr] (maxd) -- (conv);
\draw[arr] (conv) -- (norm);
\draw[arr] (norm) -- (adf);
\end{tikzpicture}
\caption{The fractional-differentiation pipeline (\texttt{DataEnricher.add\_frac\_diff\_features}). The differencing order $d$ is searched for, not assumed: three anchor assets (SPY, Gold, Crude) each yield the smallest stationarity-passing $d$ in $[0.45, 0.8]$, the maximum across anchors is adopted as a single global $d^{\ast}$, and the resulting fixed-width weight vector is applied uniformly to all eleven assets before normalization and a final per-asset stationarity check.}
\label{fig:ffd-pipeline}
\end{figure}

Figure~\ref{fig:ffd-pipeline} traces exactly this eight-step procedure end to end, from the shared log-price input through the per-anchor golden-$d$ search to the closing per-asset ADF verification. The resulting feature, \texttt{norm\_ffd\_feature}, is the only price-level feature, as opposed to a return, ratio, or oscillator, that reaches the SAA and PAA; every other feature in this section is, by design, blind to the absolute level a price series happens to be at, which is exactly the memory this feature alone is built to restore.

\subsubsection{Calendar Encoding}

A final block of six calendar features supplies periodicity without a discontinuity at period boundaries (\texttt{trigonometric\_date\_encoding}): day-of-week is encoded over a 5-trading-day cycle, day-of-month as the day's relative position within that month's actual trading days (so month length does not distort the cycle), and month-of-year over a 12-month cycle, each as a $(\sin,\cos)$ pair,

\begin{equation}
x^{\sin}_t = \sin\!\Big(\tfrac{2\pi\, k_t}{K}\Big), \qquad x^{\cos}_t = \cos\!\Big(\tfrac{2\pi\, k_t}{K}\Big),
\end{equation}

with $(k_t, K)$ equal to (day-of-week, 5), (trading-day-rank-in-month, days traded that month), and (month$-1$, 12) respectively. Unlike a raw integer day-of-week or month index, this pair-encoding places, e.g., Friday and the following Monday at a bounded cyclical distance rather than an artificial jump, and is identical in structure to the SAA's memory-block treatment of the previous action $a_{t-1}$ (Section~\ref{sec:saa}): both replace a value with a representation that does not falsely imply an unbounded numeric relationship.

Feature selection on top of this pipeline is entirely config-driven: the SAA, the PAA's per-asset token, and the PAA's portfolio token each read their own boolean feature map from the active JSON config (e.g.\ \texttt{config\_10007.json}'s \texttt{saa\_features}, \texttt{paa\_asset\_token\_features}, and \texttt{paa\_portfolio\_token\_features} sections), so architecture and feature set can be iterated on independently without touching the enrichment pipeline itself. This deliberate emphasis on feature construction, more than on the reward or the network, follows directly from the low signal-to-noise ratio of daily financial data: a well-conditioned, stationary, memory-preserving feature is worth more engineering effort than an additional backtest.

\subsection{A Shared Reward Design Methodology}
\label{sec:reward}

Both layers are trained with the same reward methodology, applied to two different notions of ``portfolio'': the SAA's private cash-plus-one-asset sub-portfolio (Section~\ref{sec:saa}), and the PAA's actual, live, eleven-asset portfolio (Section~\ref{sec:paa}). The design goal throughout is to reward risk-adjusted skill, not exposure to a rising market, and to do so with a reward that is \emph{dense}: computed from directly observable portfolio and market quantities at every single trading day, rather than only once at the end of an episode. Density matters specifically for PPO's mechanics: the algorithm's advantage estimates (Section~\ref{sec:saa}) are built from a value function trained to predict returns-to-go from every intermediate state, and a reward that stays silent for 251 of 252 steps and fires only at the very end would give that value function almost nothing to regress against until very late in training.

\subsubsection{Multi-Objective Reward Composition}

The per-step reward is a sum of five components, each isolating one aspect of trading quality, before being squashed so that no single component, or any one badly-scaled coefficient, can dominate the gradient:

\begin{equation}
R_t^{\text{raw}} = \underbrace{s_\alpha \, \Delta\ell_t}_{\text{excess return}} \;+\; \underbrace{s_\sigma \, \Delta\text{Sortino}_t}_{\text{risk-adjusted trend}} \;-\; \underbrace{\lambda_{dd}\,\Delta\text{DD}_t^{+}}_{\text{drawdown shock}} \;-\; \underbrace{\lambda_{dd}^{\text{lvl}}\,\text{DD}_t^{\text{level}}}_{\text{sustained drawdown}} \;-\; \underbrace{\lambda_{\text{exec}}\,G_t}_{\text{execution gap}}
\end{equation}

\begin{equation}
R_t = \tanh\!\Big(\frac{R_t^{\text{raw}}}{\tau}\Big) \cdot \sigma_R.
\end{equation}

The $\tanh$ squash bounds the reward PPO actually sees to $(-\sigma_R, \sigma_R)$ regardless of how large any individual raw component gets on an unusual day; $\tau$ controls how quickly the squash saturates (a larger $\tau$ keeps $R_t$ closer to linear over a wider range of $R_t^{\text{raw}}$), and $\sigma_R$ rescales the bounded output back to a magnitude PPO's default hyperparameters (clip range, value-loss coefficient) are tuned for. All five components are computed for both agents, although the two agents differ in which portfolio $\Delta\ell_t$ is measured against (Section~\ref{sec:excess-return}). Throughout development several other reward components have been tried and can be found in the code base, however if they are disabled in the best configs they are not mentioned here. 

\subsubsection{Excess Return over a Passive Counterfactual}
\label{sec:excess-return}

Rewarding raw portfolio return teaches an agent that ``the market goes up'' is a strategy, which is not the behavior this thesis is trying to elicit or measure. Both agents are instead rewarded relative to a passive counterfactual portfolio started with the same capital. Writing $V_t$ for the live (SAA: single-asset sub-) portfolio value and $B_t$ for the counterfactual's value,

\begin{equation}
\Delta\ell_t = \big(\log V_t - \log V_{t-1}\big) - \big(\log B_t - \log B_{t-1}\big).
\end{equation}

Because both terms are log-returns, $\Delta\ell_t$ is the log \emph{excess return} of the agent over ``having done nothing,'' which isolates timing and selection skill from beta exposure to a rising market. The concrete choice of counterfactual $B_t$ differs by agent: the SAA is scored against the buy-and-hold return of the one asset it happens to be trading, while the PAA is scored against a fully invested SPY buy-and-hold. Sections~\ref{sec:saa} and~\ref{sec:paa} give the exact construction of each counterfactual, since in the PAA's case there are two further diagnostic-only portfolios in the codebase that are easy to confuse with the actual reward counterfactual.

\subsubsection{Risk-Adjusted Trend: The Differential Sortino Ratio}

An exponential moving average (EMA) updates a running estimate by blending in each new observation with a fixed weight $\eta$ and discounting all older observations geometrically, $\bar x_t = \bar x_{t-1} + \eta(x_t - \bar x_{t-1})$; unlike a simple moving average over a fixed window, it needs no buffer of raw history and reacts immediately to a regime change, at the cost of a memory that fades smoothly rather than truncating sharply at a window edge. Two such EMAs track the mean step return and the downside variance (squared negative returns only, floored to avoid division blow-ups on a run that has had no losing steps yet):

\begin{equation}
\bar{r}_t = \bar{r}_{t-1} + \eta\big(r_t - \bar{r}_{t-1}\big), \qquad
\bar{v}_t^{-} = \bar{v}_{t-1}^{-} + \eta\Big(\min(r_t,0)^2 - \bar{v}_{t-1}^{-}\Big),
\end{equation}

\begin{equation}
\text{Sortino}_t = \frac{\bar{r}_t}{\sqrt{\max(\bar{v}_t^{-}, v_{\min})}}, \qquad
\Delta\text{Sortino}_t = \text{Sortino}_t - \text{Sortino}_{t-1}.
\end{equation}

The classical Sharpe ratio scores a strategy by its mean return divided by the standard deviation of \emph{all} returns, penalizing upside volatility exactly as much as downside volatility; an investor who welcomes large gains and fears only large losses is arguably better served by the Sortino ratio above, which keeps the same numerator but replaces total volatility in the denominator with \emph{downside deviation}, the standard deviation computed only over return periods that were actually negative. Rewarding the \emph{change} in this ratio, rather than its level, gives a step-wise, differentiable signal that still points the policy toward a smoother, higher-Sortino equity curve, without waiting for episode termination to score risk. A short warm-up window suppresses this term until the EMAs have enough history to be meaningful, avoiding an artificial spike from the zero-initialized baseline.

\subsubsection{Drawdown Penalties}
\label{sec:drawdown}

A \emph{drawdown} at time $t$ is the percentage decline of the current portfolio value from its own running historical peak, $\text{DD}_t = \max\big(0,\, 1 - V_t / \text{Peak}_t\big)$: the single most common risk metric in practitioner risk reporting, since, unlike volatility, it directly answers the question an investor actually asks, how much of my money was at risk at the worst point so far. Two complementary drawdown terms are used here. An incremental penalty charges only the \emph{new} drawdown incurred this step, $\Delta\text{DD}_t^{+} = \max(0, \text{DD}_t - \text{DD}_{t-1})$, over a growing lookback window (starting at 2 steps and expanding to a configured cap), so early-episode noise cannot trigger a large penalty. A persistent level penalty additionally charges every step spent underwater relative to the running peak, $\text{DD}_t^{\text{level}} = \max\big(0,\, 1 - V_t / \text{Peak}_t\big)$, discouraging the agent from lingering in a drawdown even once it has stopped worsening.

\subsubsection{The Execution Gap}
\label{sec:execgap}

The execution-gap term penalizes the distance between what the policy asked for and what the environment actually filled, net of a deadband that excludes changes too small to trade profitably; its purpose is to keep the policy honest about the cost of the requests it makes, without punishing it for a request the environment itself declined to fill because available cash or inventory was insufficient, a mechanically infeasible request rather than a strategic one. For the PAA, this is the vector gap left by the soft-execution mechanics of Section~\ref{sec:paa}:

\begin{equation}
G_t^{\text{paa}} = \sum_{j} \max\big(0,\, \lvert \Delta w_t^{(j)} - \Delta w_t^{\text{exec},(j)} \rvert - \delta\big).
\end{equation}

For the SAA, the same idea applies to the scalar action: writing $a_t^{\text{exec}}$ for the fraction of portfolio value actually traded (which can fall short of $a_t$ if cash or inventory is insufficient), the gap is

\begin{equation}
G_t^{\text{saa}} = \max\big(0,\, \lvert a_t - a_t^{\text{exec}} \rvert - g_0\big),
\end{equation}

where $g_0$ is a dynamic no-penalty band derived from the minimum trade-value threshold so that trades too small to clear the transaction-cost floor are never penalized as an execution failure. Both terms use an asymmetric multiplier, $\lambda_{\text{exec}}$ scaled up when net exposure was requested but not delivered, scaled down when the policy over-shoots and the environment simply declines to execute more than requested, without punishing either agent for the deliberate execution friction built into its own action mechanics (Sections~\ref{sec:saa} and~\ref{sec:paa}).


\subsection{Layer 1: The Single-Asset Agent (SAA)}
\label{sec:saa}

The SAA is the system's temporal specialist. Each episode it holds a position in exactly one, randomly selected asset, and its only job is to stepwise decide whether that position should grow, shrink, or stay the same, using nothing but that asset's own recent history and its own current holdings. It is the only recurrent policy in this thesis: its 2-layer Long Short-Term Memory (LSTM) cell carries a hidden state across steps within an episode, giving it a genuine, learned notion of temporal context that the feed-forward PAA (Section~\ref{sec:paa}) does not have on its own. This section gives a complete account of the SAA in isolation, exactly as it is trained, before Section~\ref{sec:paa} explains how a frozen copy of it is deployed inside the PAA.

\subsubsection{State, Observation, and Action Space}

At each step $t$ the SAA holds a position in exactly one, randomly selected asset per episode. The relevant slice of simulator state is

\begin{equation}
s_t^{\text{saa}} = \big(p_t,\, c_t,\, x_t\big),
\end{equation}

the asset's close price $p_t$, the cash balance $c_t$, and the held share count $x_t$. Writing the state this way, rather than listing everything the simulator internally tracks, isolates the \emph{endogenous} part of $\mathcal{S}$ for this agent: $c_t$ and $x_t$ are exactly the two quantities the SAA's own trades change, and $p_t$ is the one exogenous quantity its valuation and execution depend on at every step. The much larger, purely exogenous remainder of $\mathcal{S}$, the historical price and volume panel each feature of Section~\ref{sec:state} is computed from, evolves independently of the agent and is left implicit here; it still enters $\mathcal{Z}$ below, just never through anything the SAA's actions can influence. Decomposing state this way, an account state (cash, holdings) alongside a separately-described market state, is standard practice in algorithmic-trading reinforcement learning, matching, for instance, FinRL's own state definition \cite{liu2022finrldeepreinforcementlearning}. Making $s_t^{\text{saa}}$ explicit before the observation grounds the abstract POMDP tuple of Section~\ref{sec:overview} in a concrete instance, answering ``what does $\mathcal{S}$ actually contain here'' before $\mathcal{Z}$ decides what of it the network is allowed to see.

The observation $o_t^{\text{saa}} = \mathcal{Z}(s_t^{\text{saa}})$ handed to the policy is the concatenation of three blocks:


\begin{equation}
o_t^{\text{saa}} = \big[\, \mathbf{f}_t \in \mathbb{R}^{32},\ \ \mathbf{m}_t \in \mathbb{R}^{6},\ \ \mathbf{e}_i \in \{0,1\}^{11} \,\big],
\end{equation}

where $\mathbf{f}_t$ is the SAA's selected subset of the shared feature pipeline (Section~\ref{sec:state}), here the 26 market-derived indicators plus the 6 calendar features, $\mathbf{e}_i$ is a one-hot encoding of which of the $N{=}11$ universe assets is currently traded, and $\mathbf{m}_t$ is a portfolio memory block,

\begin{equation}
\mathbf{m}_t = \Big(\log\tfrac{c_t}{c_0},\ \log\tfrac{x_t p_t}{c_0},\ r_t^{\text{sub}},\ a_{t-1},\ z_t^{\text{rf}},\ \alpha_t^{\text{rf}}\Big).
\end{equation}

Each of these six entries encodes a distinct piece of information the policy needs to interpret its own recurrent state correctly. $c_0$ is the initial notional, so the first two entries are log-ratios of current cash and current asset notional against a fixed reference rather than raw dollar values, keeping their scale comparable across episodes regardless of how much capital an experiment happens to be run with. $r_t^{\text{sub}}$ is the previous-step log return of the cash-plus-asset sub-portfolio, giving the policy an immediate read on how its own last action actually performed. $a_{t-1}$ is the previous action, included because the action itself, while not part of the true simulator state $s_t^{\text{saa}}$, materially affects what a sensible next action looks like, so exposing it makes the observation behave more like a Markovian state despite $s_t^{\text{saa}}$ alone not being one. Finally, $z_t^{\text{rf}}$ is a 60-day $z$-score of the risk-free rate and $\alpha_t^{\text{rf}}$ is the sub-portfolio's excess log return over the risk-free carry. The risk-free rate here is derived from the Effective Federal Funds Rate (EFFR), the realized overnight interbank lending rate published by the Federal Reserve and the standard practitioner proxy for a nominal risk-free rate. The environment applies a daily cash-carry return derived from the EFFR (Section~\ref{sec:simulator}) so that holding cash is never modeled as a zero-return parking option, and $z_t^{\text{rf}}$, $\alpha_t^{\text{rf}}$ give the policy the same information the reward function (Section~\ref{sec:reward}) itself uses when scoring excess return, rather than forcing it to infer that context indirectly. This memory block, not the raw cash and share numbers, is what re-anchors the LSTM's hidden state to the agent's own position every step; Section~\ref{sec:saa-stability} explains why that anchoring becomes a hard requirement, not merely good practice, once the same trained weights are copied and deployed inside the PAA.

The action space is a single continuous scalar, $\mathcal{A}_{\text{saa}} = [-1, 1]$, interpreted not as a target weight but as a \emph{target position change}: a fixed fraction of portfolio notional, the action-limiting factor $\kappa_t$, is annealed over training via a three-phase linear schedule (constant at a warm-up value, linear ramp, constant at a terminal value), and $a_t = 1$ requests a buy of $1 \cdot \kappa_t \cdot c_0$ dollars, subject to available cash, while $a_t = -0.5$ requests a sell of $0.5\,\kappa_t \cdot c_0$ dollars, subject to available inventory. This design keeps the action distribution centered and bounded regardless of how large the current position already is, which matters concretely because the same scalar head, unmodified, is reused across every asset once the SAA is frozen and copied (Section~\ref{sec:paa}): a head that instead output an absolute target weight would need to already know the notional scale of whichever asset it happened to be controlling. The environment executes the implied trade, marks the sub-portfolio to market, and returns a reward that instantiates the general formula of Section~\ref{sec:reward} against a same-asset, transaction-cost-paying buy-and-hold counterfactual.

\subsubsection{Network Architecture}

\begin{figure}[htbp]
\centering
\begin{tikzpicture}[
    node distance=6mm and 8mm,
    box/.style={draw, rounded corners, minimum height=8mm, minimum width=48mm, align=center, font=\small},
    io/.style={draw, rounded corners, fill=black!5, minimum height=9mm, minimum width=48mm, align=center, font=\small},
    small/.style={draw, rounded corners, fill=black!5, minimum height=8mm, minimum width=48mm, align=center, font=\small},
    head/.style={draw, rounded corners, fill=black!10, minimum height=9mm, minimum width=44mm, align=center, font=\small},
    arr/.style={-{Latex[length=2mm]}}
]
% Two parallel input paths
\node[io] (feat) {Market + memory features\\ $[\mathbf{f}_t,\, \mathbf{m}_t] \in \mathbb{R}^{38}$\\ (32 market $+$ 6 memory)};
\node[io, right=22mm of feat] (id) {One-hot asset id\\ $\mathbf{e}_i \in \{0,1\}^{11}$};
\node[small, below=of id] (emb) {Learned embedding\\ $\text{Embedding}(11 \rightarrow 6)$};
\coordinate (topmid) at ($(feat)!0.5!(emb)$);
\node[box, below=19mm of topmid] (concat) {Concatenate\\ $[\mathbf{f}_t,\, \mathbf{m}_t,\, \text{emb}(i)] \in \mathbb{R}^{44}$};
\node[box, below=of concat] (mlp) {MLP: Linear--LayerNorm--SiLU--Dropout\\ (128 $\to$ 64), projected to $d_{f}{=}64$\\ \texttt{InputMLPFeatures}};
\node[box, below=of mlp] (lstm) {2-layer LSTM (hidden size 64)\\ \texttt{MlpLstmPolicy}};
\node[head, below left=10mm and -6mm of lstm] (actor) {Actor head\\ $\mu(a_t),\ \log\sigma(a_t)$};
\node[head, below right=10mm and -6mm of lstm] (critic) {Critic head\\ $V(s_t)$};
\node[io, below=of actor] (act) {Action $a_t \in [-1,1]$, scaled by $\kappa_t$\\ \textbf{executed by the environment}};
\node[io, below=of critic] (val) {Value estimate\\ \textbf{used only for GAE / PPO value loss}\\ never feeds back into $a_t$};
\draw[arr] (feat) -- (concat);
\draw[arr] (id) -- (emb);
\draw[arr] (emb) -- (concat);
\draw[arr] (concat) -- (mlp);
\draw[arr] (mlp) -- (lstm);
\draw[arr] (lstm) -- (actor);
\draw[arr] (lstm) -- (critic);
\draw[arr] (actor) -- (act);
\draw[arr] (critic) -- (val);
\end{tikzpicture}
\caption{SAA information flow. The observation splits into two parallel paths: raw market and memory features pass through unchanged, while the one-hot asset id is first mapped through a small learned embedding table. Both paths are concatenated before the shared MLP. A 2-layer LSTM turns the resulting feature into a recurrent state, which two independent linear heads read out: an actor head that produces the executed action, and a critic head whose value estimate is consumed only by the PPO/GAE loss and has no causal path back to the action.}
\label{fig:saa-arch}
\end{figure}

Figure~\ref{fig:saa-arch} shows the full path from observation to action, corrected to make two things explicit that a single linear pipeline would obscure. First, \texttt{InputMLPFeatures} does not feed the one-hot asset-id block into the same transformation as the market and memory features: it slices the observation into two separate tensors, sends the trailing 11-dimensional id block through a small learned embedding table (11 assets $\to$ 6 dimensions), and leaves the remaining 38-dimensional market-and-memory block untouched. Only after this split are the two tensors concatenated into a single 44-dimensional vector, which is then pushed through two dense blocks (128 then 64 units, each with LayerNorm, SiLU, and dropout) and projected to the policy's feature dimension. Second, the actor and critic heads are two independent linear readouts of the same LSTM hidden state; they do not feed into one another. The critic's value estimate $V(s_t)$ is consumed exclusively by the PPO advantage/GAE computation described below and never participates in producing $a_t$, so the two downward branches in Figure~\ref{fig:saa-arch} are deliberately drawn as terminating in two unconnected outputs.

\subsubsection{Why Actor and Critic Outputs Are Never Recombined}
\label{sec:actor-critic-theory}

A natural question, especially coming from value-based RL, is why the critic's value estimate is not added back into the actor's output to form the final action, the way a dueling network combines a state-value stream with an advantage stream into $Q(s,a) = V(s) + A(s,a) - \bar{A}(s)$. PPO, and actor-critic policy-gradient methods generally, do not follow this pattern, and the reason is more than convention.

The two heads solve two different regression problems against two different targets. The actor parameterizes a probability distribution $\pi_\theta(a_t \vert s_t)$ over actions; PPO's clipped surrogate objective is defined purely through the probability ratio $\pi_\theta(a_t\vert s_t)/\pi_{\theta_{\text{old}}}(a_t\vert s_t)$ and requires that ratio to depend only on the current policy parameters. The critic regresses $V(s_t)$ toward an empirical return-to-go target for one purpose only: producing a lower-variance advantage estimate $\hat{A}_t = \hat{R}_t - V(s_t)$ for that same surrogate. $V(s_t)$ is a training-time variance-reduction baseline, not a candidate component of the emitted action; folding it into the action would make the executed trade depend on how well the value regression happens to have converged at that point in training, an entirely different and typically noisier optimization problem than the policy's own.

Concretely, for the SAA, $V(s_t)$ is an unconstrained real number approximating a discounted sum of tanh-squashed rewards, with no reason to lie in $[-1,1]$; adding it to the actor's mean action would either require an ad hoc renormalization or would routinely push the executed target-position-change outside its valid range, injecting the critic's regression noise directly into a real trade and its transaction cost. For the PAA (Section~\ref{sec:paa-network}), the argument is sharper and purely structural rather than just a matter of numerical hygiene: the critic there produces a single scalar per state, and adding that same scalar to every one of the $N$ asset logits before the softmax of Section~\ref{sec:paa-action} would be mathematically inert, since softmax is shift-invariant, exactly the property used to justify fixing the cash logit at zero (which will be explained later). Recombining the two heads for the PAA would therefore not just be theoretically questionable; it would be a literal no-op wrapped around extra computation.

Keeping the two heads' outputs disjoint costs nothing and buys a clean separation of concerns. The actor can be evaluated with the critic discarded entirely, exactly what happens for every frozen SAA copy in Section~\ref{sec:info-flow}, which reads only the actor's mean action and never instantiates a value estimate at inference time. The critic, symmetrically, can be reweighted (via \texttt{vf\_coef}) or even pretrained on offline Monte-Carlo targets without ever perturbing the executed action. This separation does not mean the two heads are fully independent, however; Section~\ref{sec:paa-network}'s discussion of actor-critic interaction details the two indirect channels through which they still shape one another's training.

\subsubsection{Dual-Recursion Logic: LSTM Memory and Generalized Advantage Estimation}

Two pieces of PPO vocabulary the literature review introduced only in passing are worth fixing precisely before describing how the SAA's two recursions interact. The \emph{advantage} $\hat A_t = \hat R_t - V(s_t)$ compares an empirical return-to-go estimate $\hat R_t$ against the critic's own prediction $V(s_t)$ (Section~\ref{sec:actor-critic-theory}); a positive advantage means the action taken at step $t$ did better than the critic expected, and PPO's clipped surrogate objective pushes the policy to make such actions more likely going forward. \emph{Generalized Advantage Estimation} (GAE), introduced alongside PPO, computes $\hat A_t$ not from a single one-step temporal-difference residual $\delta_t = r_t + \gamma V(s_{t+1}) - V(s_t)$, but from an exponentially-weighted average of $\delta_t, \delta_{t+1}, \delta_{t+2}, \dots$ controlled by a decay parameter $\lambda$, trading off bias (small $\lambda$, leaning on a critic that may itself be miscalibrated) against variance (large $\lambda$, leaning on longer, noisier realized trajectories). \emph{Backpropagation-through-time} (BPTT) is the mechanism by which a recurrent network's gradient at step $t$ is computed by unrolling the recurrence and differentiating through every preceding step the current hidden state depends on; the longer the unroll, the more expensive the computation and the more prone it is to vanishing or exploding gradients, which is why practical recurrent training truncates it to a bounded window.

The SAA carries two nested notions of memory, and it is worth separating exactly what each one does for the credit-assignment problem, since \texttt{sb3-contrib}'s \texttt{RecurrentPPO} does not modify GAE itself. The LSTM's hidden and cell state provide within-episode memory of price dynamics: because the environment is stepped sequentially during rollout collection, with the hidden state carried forward one step at a time, the critic's value estimate $V(s_t)$ already reflects an implicit summary of everything the LSTM has seen so far in that episode, not just the current observation. GAE is then computed exactly as in standard, non-recurrent PPO: \texttt{RecurrentPPO} reuses \texttt{stable-baselines3}'s base rollout-buffer routine unmodified, running the usual backward recursion $A_t = \delta_t + (\gamma\lambda) A_{t+1}$ once over the full collected rollout (all $n_{\text{steps}}$ timesteps, across every parallel environment), masked at episode boundaries by the recorded \texttt{episode\_starts} flags. This advantage computation is never truncated or chunked; the delayed-reward attribution the differential-Sortino reward of Section~\ref{sec:reward} needs comes for free from this full-length recursion.

What \emph{is} chunked is the optimization phase, and only the optimization phase. When \texttt{RecurrentPPO} re-evaluates log-probabilities and values to compute the PPO loss, it cannot simply replay the LSTM start-to-finish over an arbitrarily long, multi-episode rollout buffer; instead it cuts each parallel environment's contiguous rollout into per-episode sequences, pads them to equal length, and replays the LSTM per sequence with its hidden state correctly reset at every recorded episode start. Backpropagation-through-time is therefore bounded to one episode's length within the current rollout. This is a deliberate and harmless decoupling: the long-horizon credit assignment lives in the advantage targets, which were already computed from a critic that saw the full sequential trajectory; the truncated BPTT during optimization only has to correct the weights that produced that already-informed value estimate, not re-derive the horizon itself.

\subsubsection{Training Dynamics and Stability}
\label{sec:saa-stability}

RecurrentPPO is sensitive to the balance between exploration and convergence speed. Learning rate, entropy coefficient, PPO clip range, and the action-limiting factor $\kappa_t$ all follow independent three-phase linear schedules (hold, ramp, hold), tuned to avoid two failure modes observed during development: early policy collapse into a long-only or short-only strategy when entropy decays too fast, and oscillatory, high-turnover behavior when $\kappa_t$ is annealed to its terminal value too early. A more structural risk is the third source of partial observability discussed in Section~\ref{sec:overview}: a temporal extractor trained on a single, isolated sub-portfolio has never seen the transition dynamics induced by a second agent's independent trades. If the frozen SAA's recurrent state were updated from the live, PAA-controlled portfolio once deployed, exogenous rebalancing by the PAA would look, from the SAA's perspective, like a domain shift it was never trained to handle. Section~\ref{sec:info-flow} shows how the shadow-portfolio design sidesteps this by construction, rather than by regularization, once a frozen copy of the SAA just described is deployed inside the PAA.

\subsection{Layer 2: The Portfolio Allocator Agent (PAA)}
\label{sec:paa}

The PAA is the system's cross-sectional specialist. Rather than learning its own opinion about each asset's timing from scratch, it consumes a frozen copy of the SAA just described, one instance per asset, as a feature, and learns how to weigh that opinion against raw market data and against every other asset's own opinion, using a self-attention mechanism (already introduced conceptually in the literature review) instead of a recurrent one. This section proceeds in the same order the PAA's information actually flows: how its observation is built from per-asset tokens (Section~\ref{sec:paa-state}), how the frozen SAA copies are integrated into that observation through a shadow-portfolio mechanism (Section~\ref{sec:info-flow}), how its action is executed (Section~\ref{sec:paa-action}), and finally how its network processes all of this end to end (Section~\ref{sec:paa-network}).

\subsubsection{State Space: Tokenizing the Asset and Portfolio Universe}
\label{sec:paa-state}

The PAA observes the full eleven-asset universe at once, rather than one asset at a time. Its observation is a sequence of $N+1$ tokens, one embedding vector per asset plus one portfolio-level token summarizing the live, actually-traded portfolio:

\begin{equation}
o_t^{\text{paa}} = \big(\mathbf{T}^{(1)}_t,\, \dots,\, \mathbf{T}^{(N)}_t,\, \mathbf{T}^{\text{pf}}_t\big), \qquad \mathbf{T}^{(i)}_t, \mathbf{T}^{\text{pf}}_t \in \mathbb{R}^{d_{\text{model}}}.
\end{equation}

Here $d_{\text{model}}$ is the shared embedding width every token is projected into before entering the attention mechanism, standard Transformer notation used by \cite{vaswani2023attentionneed} for the model's working dimensionality (here $d_{\text{model}}=64$, Section~\ref{sec:paa-network}). The two token types start out at different \emph{raw} feature counts, 29 for an asset token and 32 for the portfolio token, detailed next; only the \texttt{SAATokenizer}'s per-type linear embedding (Section~\ref{sec:paa-network}) brings both into this common $d_{\text{model}}$-dimensional space, which is what lets the shared self-attention layer compare an asset token and the portfolio token directly despite them carrying different raw information.

Each asset token $\mathbf{T}^{(i)}_t$ is built, before its final linear embedding (Section~\ref{sec:paa-network}), from that asset's 26 selected market features from the shared pipeline (Section~\ref{sec:state}), the single scalar signal $a_t^{(i)}$ the frozen SAA produces for that asset, that asset's shadow sub-portfolio holding percentage $\rho_t^{(i)}$, and the asset's current live PAA weight $w_t^{(i)}$, for a 29-dimensional raw input; $a_t^{(i)}$ and $\rho_t^{(i)}$ are exactly the two quantities the shadow-portfolio mechanism of Section~\ref{sec:info-flow} produces, and are not otherwise available anywhere else in the pipeline. The 26 market features are exactly the ones the SAA itself uses (Section~\ref{sec:saa}), minus the calendar block, and every one of them is already defined in Section~\ref{sec:state}: at each of the three horizons $N \in \{1,5,21\}$ (Section~\ref{sec:state-returns}), \texttt{log\_close\_return\_Nd}, \texttt{z\_rel\_log\_close}, \texttt{z\_vol\_breakout\_Nd}, \texttt{log\_volume\_intensity\_Nd}, and \texttt{norm\_volume\_feature\_Nd}, 15 features in total; plus eleven single-horizon indicators, \texttt{z\_log\_intra\_vol}, \texttt{norm\_closing\_strength}, \texttt{z\_log\_overnight\_gap}, and \texttt{norm\_shadow\_ratio} from the intraday bar shape, \texttt{rsi\_14}, \texttt{bb\_width\_reworked}, and \texttt{atr\_z\_norm}, the three classical technical indicators, \texttt{z\_fisher\_corr\_spy}, \texttt{z\_beta\_spy}, and \texttt{z\_market\_dispersion} from the cross-asset features, and \texttt{norm\_ffd\_feature} from the fractional-differentiation pipeline (Section~\ref{sec:state-ffd}).

The portfolio token $\mathbf{T}^{\text{pf}}_t$ is built from the six calendar features (Section~\ref{sec:state}) and 26 live-portfolio quantities, for a 32-dimensional raw input: the 12-entry live weight vector $w_t^{\text{live}}$ (cash plus 11 assets, Section~\ref{sec:paa-action}), and 14 running portfolio-risk metrics. Eight of the fourteen reuse quantities already defined earlier in this chapter, read here directly off the running episode state rather than recomputed: the current excess log return $\Delta\ell_t$ (Section~\ref{sec:excess-return}, \texttt{alpha} in the codebase), the current-step drawdown $\text{DD}_t$ (Section~\ref{sec:drawdown}), the previous- and current-step differential Sortino ratio and its underlying EMA state, $\text{Sortino}_{t-1}$, $\text{Sortino}_t$, $\bar r_t$, and $\sqrt{\bar v_t^{-}}$ (Section~\ref{sec:reward}), and the risk-free carry terms $z_t^{\text{rf}}$ and $\alpha_t^{\text{rf}}$ (Section~\ref{sec:saa}), computed here at the live-portfolio level rather than the SAA's sub-portfolio level. The remaining six are portfolio-level statistics not needed anywhere earlier in this chapter and are defined here: \texttt{sharpe}, the ratio of mean to standard deviation of portfolio returns over a trailing window, a classical risk-adjustment signal reported alongside, but computed independently of, the differential Sortino term the reward itself optimizes; \texttt{volatility}, the standard deviation of those same recent returns; \texttt{previous\_max\_drawdown}, the worst drawdown observed over a trailing window, giving the network a sense of how bad the episode has been at its worst, not only right now; \texttt{turnover}, the sum of absolute live-weight changes across assets over the recent window, a standard measure of how much rebalancing has just taken place; \texttt{concentration}, the reciprocal of the Herfindahl-Hirschman index of the live asset weights (excluding cash), normalized by the number of assets, so a value near 1 means capital is spread evenly across the universe and a value near $1/11$ means it sits almost entirely in one position; and the raw daily risk-free rate itself, alongside its already-defined $z$-scored and excess-return counterparts. The exact contents of both token types are named here in full; how they are embedded and processed thereafter is shown in simplified form in Figure~\ref{fig:paa-arch} and discussed further in Section~\ref{sec:paa-network}.

Critically, the observation $o_t^{\text{paa}}$ is recomputed from scratch every step and the PAA holds no hidden state of its own across steps: given $o_t^{\text{paa}}$, its policy is a memoryless function of the current observation alone, exactly matching the Markov assumption a standard, non-recurrent PPO update relies on. Any information about the past that reaches the PAA arrives already compressed into two places: the frozen SAA's scalar signal $a_t^{(i)}$ (Section~\ref{sec:info-flow}), and the environment's own running statistics inside the portfolio token (the Sortino EMA state, drawdown state, and so on, Section~\ref{sec:reward}). The PAA's own \texttt{AttentionEngine} (Section~\ref{sec:paa-network}) is, by contrast, a feed-forward self-attention module: it attends across the $N+1$ tokens available at the current step only, and carries nothing from one step to the next.

\subsubsection{Integrating the SAA: Shadow Portfolios}
\label{sec:info-flow}

Section~\ref{sec:saa-stability} identified a structural risk in reusing one trained SAA across eleven assets: if each copy's LSTM hidden state were updated from the \emph{live}, PAA-controlled portfolio, that hidden state would drift out of the distribution it was trained on the moment the PAA starts making allocation decisions the SAA never saw during single-asset training, an exogenous rebalance that, from the SAA's own perspective, looks like a domain shift it was never trained to handle. PyTrade avoids this risk entirely, by construction rather than by regularization, using the shadow-portfolio mechanism shown in Figure~\ref{fig:shadow-flow}.

A single frozen \texttt{RecurrentPPO} policy, the trained SAA, is evaluated once per environment step in one batched forward pass, with the $N$ assets stacked into $N$ rows; each row carries its own LSTM hidden state, indexed consistently across steps, so the batching is computationally convenient but behaviorally identical to running $N$ separate agents. Each row also owns a private, hypothetical \emph{shadow sub-portfolio}: a cash balance and share count that exist only for that row and mirror exactly what would have happened had the SAA traded that one asset alone since the start of the episode. Every step, the row's observation is rebuilt from that asset's raw market features and from memory features (Section~\ref{sec:saa}) computed against its own shadow book, never against the live, PAA-controlled portfolio. The frozen policy is then queried deterministically, and the resulting raw action, rescaled by the same action-limiting factor $\kappa$ the SAA was trained with, both (a) executes a trade against the shadow book, so it stays correctly marked to market for the next step, and (b) is emitted as the scalar signal $a_t^{(i)}$ written into asset token $i$ (Section~\ref{sec:paa-state}). Alongside it, the shadow book's own holding percentage,

\begin{equation}
\rho_t^{(i)} = \frac{\text{shadow shares}_i \cdot \text{price}_i}{\text{shadow cash}_i + \text{shadow shares}_i \cdot \text{price}_i},
\end{equation}

is computed purely from that row's isolated shadow state and injected as a second per-asset signal, distinct from the asset token's other weight field, the asset's current \emph{live} PAA weight $w_t^{(i)}$, which is read from the environment's own portfolio state, not from the shadow book. The portfolio token, by contrast, is unambiguous: it carries the actual live-traded PAA portfolio's weights and risk metrics only, and never sees any shadow-portfolio state.

Because the shadow sub-portfolios never observe the live portfolio's cash, weights, or trades, the SAA's recurrent state can never be perturbed by the PAA's own decisions: from every SAA copy's point of view, it is still playing the exact single-asset game it was trained on (Section~\ref{sec:saa}), mathematically identical to the observation distribution it saw during training. This is precisely why the memory block $\mathbf{m}_t$ of Section~\ref{sec:saa} is a required part of the SAA's own observation, and not a convenience feature: it re-anchors the LSTM's hidden state to the shadow book's own cash and position every single step, which is what keeps that hidden state observable and non-drifting once the same trained weights are frozen and copied eleven times over. The PAA, in turn, receives the SAA's opinion as an ordinary feature, alongside raw market data and the asset's live weight, and is free to trust, discount, or override it; Section~\ref{sec:paa-network}'s attention mechanism is what actually performs that weighting, and the shadow-portfolio wrapper's only job is to deliver it a clean, undistorted signal to weigh.

\begin{figure}[htbp]
\centering
\begin{tikzpicture}[
    node distance=5mm and 8mm,
    lane/.style={draw, rounded corners, fill=black!5, minimum height=8mm, minimum width=40mm, align=center, font=\footnotesize},
    saabox/.style={draw, rounded corners, fill=black!15, minimum height=8mm, minimum width=40mm, align=center, font=\footnotesize},
    arr/.style={-{Latex[length=1.8mm]}},
    every node/.style={font=\footnotesize}
]
% Lane 1
\node (l1top) {Asset 1};
\node[lane, below=of l1top] (shadow1) {Shadow sub-portfolio 1\\ cash$_1$, shares$_1$ (isolated)};
\node[lane, below=of shadow1] (obs1) {Rebuild SAA obs.\\ market feats. + memory$_1$ + id$_1$};
\node[saabox, below=of obs1] (saa1) {Frozen SAA\\ LSTM state row $r{=}1$};
\node[lane, below=of saa1] (sig1) {Signal $a_t^{(1)} = \kappa \cdot \hat{a}^{(1)}$\\ Shadow weight $\rho_t^{(1)}$};
% Lane i
\node[right=18mm of l1top] (litop) {Asset $i$ ($\cdots$)};
\node[lane, below=of litop] (shadowi) {Shadow sub-portfolio $i$\\ cash$_i$, shares$_i$ (isolated)};
\node[lane, below=of shadowi] (obsi) {Rebuild SAA obs.\\ market feats. + memory$_i$ + id$_i$};
\node[saabox, below=of obsi] (saai) {Frozen SAA (shared weights)\\ LSTM state row $r{=}i$};
\node[lane, below=of saai] (sigi) {Signal $a_t^{(i)} = \kappa \cdot \hat{a}^{(i)}$\\ Shadow weight $\rho_t^{(i)}$};
% Lane N
\node[right=18mm of litop] (lNtop) {Asset $11$};
\node[lane, below=of lNtop] (shadowN) {Shadow sub-portfolio 11\\ cash$_{11}$, shares$_{11}$ (isolated)};
\node[lane, below=of shadowN] (obsN) {Rebuild SAA obs.\\ market feats. + memory$_{11}$ + id$_{11}$};
\node[saabox, below=of obsN] (saaN) {Frozen SAA (shared weights)\\ LSTM state row $r{=}11$};
\node[lane, below=of saaN] (sigN) {Signal $a_t^{(11)} = \kappa \cdot \hat{a}^{(11)}$\\ Shadow weight $\rho_t^{(11)}$};
% PAA box spanning all lanes
\node[draw, rounded corners, fill=black!10, minimum width=132mm, minimum height=12mm, below=10mm of sigi, align=center] (paa) {PAA \texttt{SAATokenizer} injects $a_t^{(i)}$ and $\rho_t^{(i)}$ into asset token $i$ $\to$ \texttt{AttentionEngine} $\to$ live portfolio weights $\mathbf{w}_t$};
\draw[arr] (shadow1) -- (obs1);
\draw[arr] (obs1) -- (saa1);
\draw[arr] (saa1) -- (sig1);
\draw[arr] (sig1.south) .. controls +(-23mm,-16mm) and +(-6mm,6mm) .. (shadow1.west |- sig1) node[midway, left, xshift=-5mm, yshift=9mm] {commit shadow trade};
\draw[arr] (shadowi) -- (obsi);
\draw[arr] (obsi) -- (saai);
\draw[arr] (saai) -- (sigi);
\draw[arr] (shadowN) -- (obsN);
\draw[arr] (obsN) -- (saaN);
\draw[arr] (saaN) -- (sigN);
\draw[arr] (sig1) -- (paa.north -| sig1);
\draw[arr] (sigi) -- (paa.north -| sigi);
\draw[arr] (sigN) -- (paa.north -| sigN);
\end{tikzpicture}
\caption{Shadow-portfolio signal generation. One frozen SAA (shared weights, one LSTM state per asset row) trades an isolated, hypothetical sub-portfolio for every asset in parallel. Only the resulting scalar action and the shadow book's own holding percentage, never the shadow book itself, cross into the PAA's asset tokens; the live, PAA-controlled portfolio never touches any SAA's recurrent state.}
\label{fig:shadow-flow}
\end{figure}

\subsubsection{Action Space: Softmax Allocation with a Fixed Cash Anchor}
\label{sec:paa-action}

The PAA's actor head outputs $N$ real-valued logits $\mathbf{z}_t \in \mathbb{R}^N$, one per risky (non-cash) asset, $\mathcal{A}_{\text{paa}} = \mathbb{R}^N$. A cash logit is fixed at $z_t^{(0)} \equiv 0$ and prepended, and the executed target weights follow from a stable softmax, the standard mapping from an unconstrained real-valued vector onto the probability simplex (all entries non-negative, summing to one):

\begin{equation}
w_t^{(j)} = \frac{\exp\big(z_t^{(j)} - \max_k z_t^{(k)}\big)}{\sum_{k=0}^{N} \exp\big(z_t^{(k)} - \max_k z_t^{(k)}\big)}, \qquad j = 0, \dots, N.
\end{equation}

Anchoring the cash logit at a fixed value, rather than learning it, removes one redundant degree of freedom from the policy: softmax is shift-invariant, $\text{softmax}(\mathbf{z}) = \text{softmax}(\mathbf{z}+c)$ for any constant $c$, so an unconstrained, learned cash logit would only ever matter relative to the asset logits and could always be re-absorbed into them, meaning the network would have to learn to keep it at some implicit reference value anyway. Fixing it at zero removes that redundancy outright and gives the agent a well-defined, always-available ``do nothing'' reference point that does not have to be rediscovered from scratch at every stage of training.

The resulting target weights are not applied instantaneously. A soft-execution step moves the live portfolio only part-way toward the target,

\begin{equation}
\Delta w_t = w_t - w_t^{\text{live}}, \qquad
\Delta w_t^{\text{exec}} = \eta_{\text{exec}} \cdot \Delta w_t \cdot \mathbb{1}\big[\lvert \Delta w_t \rvert > \delta\big],
\end{equation}

with step size $\eta_{\text{exec}} \in (0,1]$ and deadband $\delta$. The deadband suppresses trades too small to be worth the transaction cost (Section~\ref{sec:friction}), and the partial step spreads a large rebalance over several days rather than realizing it as one trade, which is closer to how a real allocator would rebalance and keeps single-step turnover, and therefore cost, bounded. The reward for this action instantiates the general formula of Section~\ref{sec:reward} against a fully invested SPY buy-and-hold counterfactual, using the vector form $G_t^{\text{paa}}$ of the execution-gap term (Section~\ref{sec:execgap}), since $\Delta w_t$ and $\Delta w_t^{\text{exec}}$ above are exactly the raw and realized weight deltas that formula compares.

One further clarification matters here because it is easy to conflate three similarly-named portfolios in the codebase. The PAA's actual excess-return counterfactual, $B_t$ in Section~\ref{sec:excess-return}, is a $100\%$-invested, transaction-cost-paying SPY buy-and-hold: an SPY-only position sized with the full initial portfolio value, which pays its own one-time transaction cost at episode reset and is then never traded again, only marked to market every step. A separate \emph{comparison portfolio} is a static clone that receives the exact same initial cash and positions as the live portfolio at episode reset (including whatever random or cash-only starting allocation was sampled for that episode) and is then never traded again; it is reported for diagnostics (\texttt{comparison\_final\_value}) but does not enter the reward. A further, fixed-weight \emph{benchmark portfolio} ($45\%$ SPY, $20\%$ Gold, $10\%$ EWJ, $10\%$ EWG, $5\%$ Crude, $5\%$ EWQ, $5\%$ EWT, fully invested from day one) is also simulated every episode and reported as a diagnostic ``alpha'' metric, but likewise never enters the reward. The PAA is trained to beat the SPY buy-and-hold specifically, not either of these two diagnostic-only portfolios, which is consistent with the SAA's own same-asset buy-and-hold counterfactual: in both cases, $B_t$ is the simplest passive thing an investor could actually do with this capital.

\subsubsection{Network Architecture: Shared Attention with Private Heads}
\label{sec:paa-network}

\begin{figure}[htbp]
\centering
\begin{tikzpicture}[
    node distance=7.2mm and 7.2mm,
    tok/.style={draw, rounded corners, fill=black!5, minimum height=9mm, minimum width=30mm, align=left, font=\scriptsize, inner sep=2mm},
    box/.style={draw, rounded corners, minimum height=8mm, minimum width=96mm, align=center, font=\small},
    head/.style={draw, rounded corners, fill=black!10, minimum height=9mm, minimum width=34mm, align=center, font=\small},
    arr/.style={-{Latex[length=2mm]}}
]
\node[tok] (t1) {\textbf{Asset token 1}\\[1mm]
26 market features\\ (Section~\ref{sec:paa-state})\\[1mm]
$+$ SAA signal $a_t^{(1)}$\\
$+$ shadow holding \% $\rho_t^{(1)}$\\
$+$ live PAA weight $w_t^{(1)}$};
\node[tok, right=of t1] (t2) {\textbf{Asset token $i$}\\[1mm] $(\cdots$ identical 26 features\\ $+\ a_t^{(i)} + \rho_t^{(i)} + w_t^{(i)})$};
\node[tok, right=of t2] (t3) {\textbf{Asset token 11}\\[1mm] $(\cdots$ identical 26 features\\ $+\ a_t^{(11)} + \rho_t^{(11)} + w_t^{(11)})$};
\node[tok, right=of t3] (tp) {\textbf{Portfolio token}\\[1mm]
6 calendar features\\[1mm]
$+$ 12 live weights (cash $+$ 11)\\
$+$ 14 portfolio-risk metrics\\ (Section~\ref{sec:paa-state})};
\node[fit=(t1)(t2)(t3)(tp), inner sep=0pt] (tokrow) {};
\node[box, below=13mm of tokrow] (embed) {\texttt{SAATokenizer}: Linear embed each token to $d_{\text{model}}{=}64$\\ + scaled learned asset-id embedding + LayerNorm};
\node[box, below=of embed] (shared) {Shared \texttt{TransformerEncoder} (2 layers, 4 heads, Pre-LN, GELU)\\ self-attention across all $N{+}1$ tokens};
\coordinate (branchrow) at ($(shared.south)+(0,-12mm)$);
\node[head, anchor=west] (actorL) at (t1.east |- branchrow) {Private actor layer\\ (asset tokens only) $\to$ mean-pool};
\node[head, anchor=east] (criticL) at (tp.west |- branchrow) {Private critic layer\\ (portfolio token only)};
\node[head, below=of actorL] (actorH) {Actor head: $N$ logits $\mathbf{z}_t$\\ + fixed cash logit $0$ $\to$ softmax $\mathbf{w}_t$};
\node[head, below=of criticL] (criticH) {Critic head: $V(s_t)$};
\draw[arr] (t1) -- (embed);
\draw[arr] (t2) -- (embed);
\draw[arr] (t3) -- (embed);
\draw[arr] (tp) -- (embed);
\draw[arr] (embed) -- (shared);
\draw[arr] (shared) -- (actorL);
\draw[arr] (shared) -- (criticL);
\draw[arr] (actorL) -- (actorH);
\draw[arr] (criticL) -- (criticH);
\end{tikzpicture}
\caption{PAA information flow. Every asset token carries the same 26 market features (named in full in Section~\ref{sec:paa-state}), the injected SAA signal, the shadow sub-portfolio's own holding percentage, and the asset's current live PAA weight; the portfolio token carries 6 calendar features, the full live weight vector, and 14 running portfolio-risk metrics (also named in Section~\ref{sec:paa-state}). All tokens are embedded to equal width, attend to one another in a shared transformer encoder, then diverge into private single-layer encoders, horizontally aligned with the asset/portfolio token span, that isolate actor and critic gradients before the final linear heads.}
\label{fig:paa-arch}
\end{figure}

Figure~\ref{fig:paa-arch} shows the PAA's \texttt{AttentionEngine}; the exact feature names behind each token were given above (Section~\ref{sec:paa-state}) and are configured in \texttt{paa\_asset\_token\_features} and \texttt{paa\_portfolio\_token\_features}. A learned per-asset identity embedding, scaled by a small trainable gate and added before a LayerNorm, lets the network distinguish assets without the identity signal overwhelming the market-derived content of the token, an issue that surfaced early in development when an unscaled identity embedding dominated the token representation. All $N+1$ tokens then pass through a shared, two-layer \texttt{TransformerEncoder} so that every asset's representation can be conditioned on every other asset and on the live portfolio's own state, i.e., cross-sectional attention.

The architecture then deliberately splits. Rather than reading the actor and critic heads directly off the shared encoder's output, each head gets its own single \texttt{TransformerEncoderLayer} that only ever receives gradient from its own loss. The actor's private layer only sees the $N$ asset tokens (mean-pooled afterwards into the policy's latent vector); the critic's private layer only sees the portfolio token. This absorbs the well-known tendency of actor and critic losses to pull a shared trunk in conflicting directions before that conflict reaches the (much more expensive to retrain) shared attention stack. The actor head's $N$ raw logits are combined with the fixed cash anchor and passed through the softmax of Section~\ref{sec:paa-action} to produce the executed target weights; the critic head reads only the portfolio token into a single scalar value estimate.

\subsubsection{Actor-Critic Interaction}

As formalized in Section~\ref{sec:actor-critic-theory}, adding the critic's scalar value estimate to every asset logit before the softmax of Section~\ref{sec:paa-action} would change nothing about the resulting weights at all, since softmax is shift-invariant; recombining the heads here would therefore be mathematically inert, not merely unconventional. There is, deliberately, no direct connection between the actor and critic heads: neither reads the other's output, and at inference time the value estimate has no causal path to the executed action, mirroring the SAA's actor/critic separation in Figure~\ref{fig:saa-arch}. The two heads interact only through two indirect channels. First, both private encoder layers are downstream of the same shared, two-layer \texttt{TransformerEncoder}; gradients from the policy loss and from the value loss both flow into that shared trunk's weights, so a representation collapse driven by one loss can, in principle, degrade the other. This is precisely the risk the private per-head layers of this section are designed to absorb: they let the two losses disagree at the last layer without forcing that disagreement all the way back into the expensive shared attention stack every update. Second, the two losses are combined into a single scalar by PPO's objective, $L = L_{\text{policy}} - c_1 L_{\text{value}} + c_2 L_{\text{entropy}}$, so the critic still shapes the actor's training indirectly through the advantage estimates $\hat{A}_t = R_t - V(s_t)$ used inside $L_{\text{policy}}$: a miscalibrated critic yields a noisy or biased advantage signal, which is the standard way actor-critic methods couple even when their network paths are otherwise disjoint. Xavier initialization with GPT-2-style residual-branch scaling is applied specifically to the shared trunk to keep this coupling numerically stable as depth increases.

\subsection{Simulator Design}
\label{sec:simulator}

Every execution mode described below runs on top of one custom Gymnasium environment, so it is worth fixing what that means before describing the modes themselves. Gymnasium, the maintained successor to OpenAI Gym, is the standard interface reinforcement-learning environments implement. It exposes two core methods: \texttt{reset()}, which starts a new episode and returns the first observation, and \texttt{step(action)}, which advances the environment by one step and returns the next observation, a reward, two termination flags, and a diagnostic info dictionary. This interface is deliberately framework-agnostic: it exchanges plain NumPy arrays, not PyTorch tensors, so any policy library can wrap it without the environment itself depending on a deep-learning framework. Stable-Baselines3 and sb3-contrib, the PyTorch-based PPO and RecurrentPPO implementations used throughout this thesis (Sections~\ref{sec:saa} and~\ref{sec:paa}), convert those arrays to and from PyTorch tensors internally; the environment itself never has to. 

Off-the-shelf RL trading frameworks typically expose a single Gym-style environment with a fixed observation interface and a fixed action parameterization. In particular, FinRL’s portfolio-optimization setting commonly represents the action as continuous portfolio weights (e.g., implemented via \texttt{gym.spaces.Box}) and executes trades accordingly \cite{liu2022finrldeepreinforcementlearning}. This design is well aligned with standard single-policy training loops, but it does not directly work with the specialized requirements of the hierarchical setup. In this case, the environment must support several \emph{modes} under a shared implementation. Concretely, it must (i) switch by configuration between a single-asset-agent (SAA) formulation and a portfolio-allocator-agent (PAA) formulation, (ii) maintain \emph{shadow} (hypothetical) states and sub-portfolios for internal rollouts or private accounting, and (iii) allow interactions with multiple agents, including frozen sub-policies, while reusing a single feature cache and a single transaction-cost model for consistency across modes. Supporting these capabilities within one testable interface is a significant specialization compared to typical off-the-shelf designs, even though such behavior could in principle be implemented by extending existing frameworks with non-trivial engineering effort. Building \texttt{TradingEnv} from scratch is what makes that hierarchical setup possible within one consistent, testable codebase, and the rest of this section describes exactly how.

\subsubsection{Execution Modes and Data Infrastructure}

\texttt{TradingEnv} supports four execution modes, selected by configuration rather than by subclassing: a single-asset target-position mode used to train the SAA, a portfolio-weights mode used to train the PAA, and two instruction-list modes (simple and tranche execution) reserved for manual or scripted trading experiments outside of the RL loop. They were used to test rule-based trading algorithms to verify the execution logic of the environment. A \texttt{MarketDataCache} precomputes and memory-maps the entire feature panel once per run, so that repeated episode resets during training never re-run feature engineering; an \texttt{EpisodeBuffer} records every step's prices, weights, and portfolio metrics for the reward and observation functions to read back a rolling window from, without recomputation. Every episode is sampled from within a single contiguous date block; how those blocks are carved out of the full history and kept separate from one another is the subject of the next subsection.

\subsubsection{Purged, Block-Level Train/Validation/Test Splitting}

The split itself is constructed once, by \texttt{MarketDataCache.\_create\_time\_series\_blocks}, directly on the full, continuous date range, before any episode is ever sampled. It is not a single chronological cut into three pieces; it repeats a purged train/validation cycle,

\begin{equation}
[\text{train}_0][\text{purge}][\text{val}_0]\,[\text{train}_1][\text{purge}][\text{val}_1]\,\cdots\,[\text{purge}][\text{test}],
\end{equation}

across the available history, with a configurable target validation share (here $20\%$ of every train\,+\,validation super-block, \texttt{train\_val\_split\_ratio}) and a fixed-length purge gap ($60$ trading days, \texttt{purge\_length\_days}) inserted between each training block and the validation block that follows it. Each super block's length is (within the limits of integer division) similar. A single out-of-sample test block is then held out at the very end of the date range, itself preceded by one more purge gap. Block boundaries, not individual days, are the unit of assignment: every training, validation, and test episode is sampled (\texttt{sample\_episode\_start}, weighted by each block's day count from entirely within one block, and an episode's full length (252 trading days, \texttt{episode\_length\_days}) is always contained inside that block's date range, never straddling a purge gap or a block boundary.

The purge gap is not a cosmetic buffer; it is sized against the same rolling-window features described in Sections~\ref{sec:state-returns} and~\ref{sec:state-ffd}. Every engineered feature is computed once, causally, over the full continuous time series, using only backward-looking windows of up to $w_{\text{norm}}=36$ days (and, for the fractionally differentiated feature, up to $\ell^{\ast}\le 120$ days). Without a gap at least as wide as the longest such window, the first days of a validation block would carry features whose rolling statistics were computed partly from data inside the immediately preceding training block, letting information the model was fitted on leak into the observations it is being evaluated on, exactly the mechanism López de Prado \cite{LopezdePrado.2018} identifies as the reason ordinary $k$-fold cross-validation fails on financial time series: samples are serially correlated and overlap through shared lookback windows, so a random, non-contiguous split systematically overstates out-of-sample performance. The $60$-day purge used here exceeds the dominant $36$-day feature window with margin to spare; the FFD feature's rarely-approached $120$-day cap is the one place this margin is not formally guaranteed, which is worth flagging as an assumption rather than a proof. No purge is inserted on the reverse edge, between a validation block and the following training block, because the risk purging defends against is asymmetric: it is validation (and test) performance, not training loss, that this thesis reports and selects checkpoints on, so only training-into-evaluation leakage needs to be closed off.

This design follows López de Prado's broader prescription for backtesting DRL and other sequential trading systems on a purpose-built simulator \cite{LopezdePrado.2018}: chronological, purged, block-level splitting in place of shuffled cross-validation; several validation blocks spread across the full history rather than one arbitrary cutoff, so that validation performance reflects multiple market regimes (e.g., the dot-com decline, the 2008 crisis, and the 2020 crash all fall inside different blocks of the 2000--2025 range used here) instead of a single one; and a final test block that is touched exactly once. Because checkpoint selection never touches that final block, as just noted, the reported test-set result is a single, previously-unseen evaluation rather than the best of several attempts against the same data, the specific failure mode López de Prado terms backtest overfitting. This implementation uses repeated purged train/validation cycles rather than his full Combinatorial Purged Cross-Validation (CPCV): CPCV's combinatorial fold count would make an already expensive PPO training run intractable to repeat many times over, and repeatedly re-running full training against many combinatorial folds would itself reintroduce a version of the same overfitting risk it is meant to prevent. The scheme adopted here is best read as the practical, purged walk-forward compromise López de Prado himself allows for settings where full CPCV is computationally out of reach.

\subsubsection{Hardware Placement: CPU Simulation, GPU Inference}

\begin{figure}[htbp]
\centering
\begin{tikzpicture}[
    node distance=6mm and 10mm,
    cpu/.style={draw, rounded corners, fill=black!5, minimum height=9mm, minimum width=48mm, align=center, font=\small},
    gpu/.style={draw, rounded corners, fill=black!12, minimum height=9mm, minimum width=48mm, align=center, font=\small},
    grp/.style={draw, dashed, rounded corners, inner sep=3mm},
    arr/.style={-{Latex[length=2mm]}}
]
\node[cpu] (cache) {\texttt{MarketDataCache}\\ (feature panel, once per run)};
\node[cpu, below=of cache] (envs) {$k$ $\times$ \texttt{TradingEnv} in \texttt{SubprocVecEnv}\\ (one OS process per worker)};
\node[cpu, below=of envs] (norm) {\texttt{VecNormalize}\\ (running obs/reward statistics)};
\node[gpu, below=of norm] (wrapper) {\texttt{SAASignalWrapper}\\ batches obs across all (env, asset) rows};
\node[gpu, below=of wrapper] (saa) {Frozen SAA policy on GPU\\ one batched forward pass, $k{\times}N$ rows};
\node[gpu, below=of saa] (ppo) {PAA \texttt{AttentionEngine} (PPO policy) on GPU\\ batched forward/backward across $k$ envs};
\node[cpu, below=of ppo] (step) {\texttt{env.step()} back in each worker process\\ execution, cost model, reward, bookkeeping};
\draw[arr] (cache) -- (envs);
\draw[arr] (envs) -- (norm);
\draw[arr] (norm) -- (wrapper);
\draw[arr] (wrapper) -- (saa);
\draw[arr] (saa) -- (ppo);
\draw[arr] (ppo) -- (step);
\draw[arr] (step.east) .. controls +(30mm,0) and +(30mm,0) .. (envs.east);
\begin{scope}[on background layer]
\node[grp, fit=(cache)(envs)(norm)(step), label={[font=\small]above:CPU: many cheap, independent workers}] {};
\node[grp, fit=(wrapper)(saa)(ppo), label={[font=\small]right:GPU: few, large, regular tensor ops}] {};
\end{scope}
\end{tikzpicture}
\caption{Class interaction and hardware placement of the simulator/training loop for PAA training (equivalent to a high-level component diagram). Environment simulation is deliberately kept on CPU workers; every neural forward/backward pass is batched across environments (and, for the frozen SAA, across assets) into large GPU tensor operations before control returns to the CPU-side environments.}
\label{fig:sim-classes}
\end{figure}

Figure~\ref{fig:sim-classes} makes explicit why this split between CPU and GPU work matters on a CUDA-enabled machine, beyond the generic "neural networks like GPUs" argument. Environment logic (trade execution, the cost model of Section~\ref{sec:friction}, reward bookkeeping) is scalar, branch-heavy, and cheap; running it on a GPU would waste the device's parallelism on work it is not shaped for, so each \texttt{SubprocVecEnv} worker keeps this logic on its own CPU process, and $k$ workers advance $k$ independent episodes concurrently. Every neural network evaluation, by contrast, is deliberately batched before it touches the GPU: \texttt{SAASignalWrapper} stacks the $k$ parallel environments' $N$ assets into a single $(k \cdot N)$-row tensor and queries the frozen SAA once per environment step, rather than once per asset; the PAA's own \texttt{AttentionEngine} likewise processes all $k$ environments' token sequences in one batched forward (and, during optimization, backward) pass. This turns what would otherwise be $k \cdot (N{+}1)$ small, latency-bound calls per step into two large, regular matrix multiplications, which is exactly the workload profile a GPU is built for. The net effect is that the expensive part of the pipeline (two neural networks, one of them evaluated per-asset) is amortized over a single batched CUDA call per step, while the cheap part (environment transition dynamics) stays where it is fast: on ordinary CPU cores, in parallel, never blocking on the GPU except to hand off and receive back a batch of observations and actions.

\subsection{Friction Modeling}
\label{sec:friction}

Every executed trade pays three cost components, all proportional to traded notional, plus an optional fixed per-order fee:

\begin{equation}
\text{TC}_t = \underbrace{c_{\text{comm}} \sum_j n_j}_{\text{commission}} + \underbrace{c_{\text{spread}} \sum_j n_j}_{\text{half-spread}} + \underbrace{\lambda_{\text{impact}} \sum_j n_j \sqrt{\tfrac{n_j}{\text{ADV}_j}}}_{\text{market impact}} + f_{\text{fixed}} \cdot \#\{\text{orders}\},
\end{equation}

where $n_j$ is the traded notional for asset $j$ and $\text{ADV}_j$ is that asset's trailing \emph{average daily (dollar) volume} over a configurable window, the standard liquidity measure practitioners use to judge how large a trade is relative to what a market can typically absorb in one session. The commission and half-spread terms are the most familiar frictions: a broker fee and half the quoted bid-ask spread, both linear in traded notional. The market-impact term captures a different, size-dependent effect: a sufficiently large trade moves the price it gets filled at, and this cost is modeled, following standard market-microstructure practice, as growing with the \emph{square root} of trade size relative to available liquidity, $n_j\sqrt{n_j/\text{ADV}_j}$, rather than linearly; empirical microstructure studies consistently find impact to be sublinear in trade size, so a trade twice as large costs less than twice as much per dollar traded, since only the marginal, unusually large portion of it has to walk further up (or down) the order book. Commission, half-spread rates, the impact and fixed per-order coefficients
are independently configurable per experiment; setting all four to zero recovers a frictionless environment useful for isolating whether a policy change affects raw signal quality or only realized net returns after costs.

\subsection{Asset Universe}
\label{sec:universe}

The current universe spans 11 liquid, exchange-traded instruments covering equities, a commodity complex, and precious metals: SPY (US equities), EWG, EWH, EWJ, EWQ, EWS, EWT, EWU, and EWY (Germany, Hong Kong, Japan, France, Singapore, Taiwan, the UK, and South Korea equities, respectively), Crude oil, and Gold. Every one of these eleven instruments is a USD-denominated, US-exchange-listed vehicle: the country ETFs are iShares MSCI products quoted and settled in dollars on US exchanges, not local-currency instruments on their respective home markets, and Crude and Gold are likewise tracked through USD-quoted contracts. The simulator's data pipeline includes a general-purpose FX-conversion step precisely so that a future, more geographically diverse universe could mix local-currency instruments, but for this concrete 11-asset universe every entry in the currency map already resolves to USD, so that conversion step is a documented no-op here rather than an active part of the pipeline.

This choice of universe serves the thesis's research question directly. Restricting the universe to already-USD-denominated instruments removes currency risk as a confounding factor: any cross-asset structure the PAA's attention mechanism discovers can be attributed to genuine co-movement in the underlying assets, not to an artifact of an FX conversion the model never has to reason about. All eleven instruments are ETFs or exchange-traded futures/ETCs rather than single stocks, which avoids single-name idiosyncratic risk and corporate-action noise (splits, spin-offs, individual earnings surprises) that a single-stock universe would otherwise force the feature-engineering layer to absorb. Each instrument also has a long, continuously liquid daily history spanning the 2000--2025 window used throughout this thesis, which is a precondition for the rolling-window statistics of Section~\ref{sec:state} to be well-estimated from the earliest usable episode onward. Finally, the eleven instruments were chosen for cross-sectional \emph{diversity}, not for any expected relative-value relationship between specific pairs: equities from seven distinct regions, a broad commodity, and a traditional safe-haven metal give the PAA a universe in which regime-dependent co-movement (e.g., regional equity indices moving together while gold diverges in risk-off periods) is plausible but not hand-engineered into the reward or the features. A genuinely useful cross-sectional allocator should be able to discover that structure through attention, rather than have it handed to it.

